\documentclass[10pt,a4paper]{amsart}
\usepackage[margin=19mm]{geometry}
\usepackage{amsmath,amssymb,mathtools}
\usepackage[scr=rsfs,cal=euler]{mathalfa}
\usepackage{xfrac}
\usepackage[unicode,hidelinks,bookmarks=false]{hyperref}
\newcommand{\ie}{i.e.\ }
\newcommand{\cf}{cf.\ }
\newcommand{\resp}{respectively\ }
\newcommand{\Subsection}{\subsection}
\providecommand{\nobreakdash}{\nobreak\hbox{-}}
\setlength{\emergencystretch}{2em}
\multlinegap0pt
\usepackage{amssymb}
\usepackage{bm}
\usepackage{xcolor}
\usepackage[shortlabels]{enumitem}
\usepackage{mathtools}
\newtheorem{lemma}{Lemma}
\DeclareMathOperator{\crys}{crys}
\DeclareMathOperator{\homm}{Hom}
\title{On the $v$-adic values of G-functions II:\\ \Large Towards Effective Brauer-Siegel}
\author{Georgios Papas}
\date{Source: arXiv:2510.11814v2 (CC BY 4.0)}
\begin{document}
\maketitle

\noindent\emph{Source and licence.} On the $v$-adic values of G-functions II:\\ \Large Towards Effective Brauer-Siegel. Version arXiv:2510.11814v2 (CC BY 4.0), distributed by arXiv under the Creative Commons Attribution 4.0 International licence, http://creativecommons.org/licenses/by/4.0/, as recorded in the arXiv licence metadata for this exact version and shown on the abstract page https://arxiv.org/abs/2510.11814v2. Full licence notice: \texttt{licenses/arxiv-2510.11814-CC-BY-4.0.txt}.\par\smallskip

\noindent\textit{Editorial scope.} Counted unit: the article's lemma lemmacrysordbasis (source0.tex lines 292--295). The article's complete original proof of it is delivered with it (source0.tex lines 296--299, one \texttt{proof} environment of 4 lines). No mathematics is written, repaired or re-derived: the statement and the proof are the article's own lines, in the article's own notation, and no retained line is substituted or re-flowed.  No further source-local context is needed: every cross-reference of every retained line resolves inside the two delivered spans.  Nothing was omitted from any retained span, so the package carries no omission record.  No retained line contains a citation, so the wrapper carries no bibliography and no bibliography entry.  No retained line contains a figure environment, an image-inclusion command or drawing code, so no image is delivered and the package declares no asset dependency.  Editorial formatting only, in addition: an AMS \texttt{amsart} preamble that restates the article's own declarations and macros used by the retained lines, namely the environments \texttt{lemma} on separate counters, together with the standard packages those lines need.  The retained environments print in this document as Lemma 1 in order of appearance, whereas the article prints the same items with the numbering of its own full document.  The complete native source of the article as distributed in the arXiv e-print for this exact version is bundled unchanged as \texttt{sources/186990/source0.tex}.
\begin{lemma}\label{lemmacrysordbasis}Let $\tilde{E}$ and $\tilde{E}'$ be ordinary elliptic curves over $k_v$ and let $f\in \homm(\tilde{E},\tilde{E}')$ be a non-zero homomorphism. Let $\Gamma_v(\tilde{E})=\{\gamma,\delta\}$ and $\Gamma_v(\tilde{E}')=\{\gamma',\delta'\}$ be chosen as above. 
	
	Then, for its pullback $f_{\crys}:H^1_{\crys}(\tilde{E}'/K_{v,0})\rightarrow H^1_{\crys}(\tilde{E}/K_{v,0})$, there exist $\zeta_0$, $\zeta_1\in K_{v,0}$ such that $f_{\crys}(\gamma')=\zeta_1\cdot \gamma$ and $f_{\crys}(\delta')=\zeta_0\cdot \delta$.
\end{lemma}

\begin{proof} The pullback $f_{\crys}$ of $f$ on the level of crystalline cohomology will be an endomorphism in the category of $F$-isocrystals. In other words, we have $f_{\crys}\circ\phi'_{\crys} =\phi_{\crys}\circ f_{\crys}$, where $\phi'_{\crys}$ stands for the $q$-th power Frobenius of $H^1_{\crys}(\tilde{E}'/K_{v,0})$.
	
	 Applying this to $\gamma'$ it is easy to see that we must have that $\phi_{\crys}( f_{\crys}(\gamma'))= \lambda_1\cdot f_{\crys}(\gamma')$. Therefore, since $\phi_{\crys}$ is linear with distinct eigenvalues, we get that there exists $\zeta_1\in K_0$ such that $f_{\crys}(\gamma^{\prime})=\zeta_1\cdot \gamma$. The existence of $\zeta_0$ for $\delta$ follows similarly.
\end{proof}

\end{document}
